\input{preamble.tex}
\title{Orbit intersections and marked semidirect products}
\author{Human Compression Project}\date{}
\begin{document}\maketitle
\begin{abstract}
Orbit intersections give a common construction of flag-transitive geometries
for semidirect products, including actions with a nontrivial chamber
stabilizer. We prove all parabolic and panel formulas, residual connectivity
and the faithful regular specialization. Element orbits and orbital Coxeter
presentations supply the diagonal constructions. For signed permutation
markings, a separate nested-module argument and an explicit cocycle change
of coordinates recover six families, including the odd-rank case left
unverified in the source. This edition extends the complete earlier regular
theory; it remains a private draft pending source and proof review.
\end{abstract}
\section{Conventions and the two intersection properties}
All type sets in this paper are finite. Cosets are left cosets. In a coset geometry $(B,(B_i)_{i\in I})$ write $B_J=\bigcap_{i\in J}B_i$, with $B_\varnothing=B$. A regular hypertope is thin, residually connected and flag-transitive. Automorphisms preserve types. In its faithful standard model, the group acts freely and transitively on chambers. Its distinguished involutions $b_i$ interchange a fixed chamber and its $i$-adjacent chamber; then
\begin{equation}\label{eq:parabolic}
B_J=\langle b_i:i\in I\setminus J\rangle,\qquad B_I=1.
\end{equation}
We recall why the geometric hypotheses give these assertions. In a thin geometry a type-preserving automorphism fixing a chamber fixes all its adjacent chambers. Residual connectivity implies that the chamber graph, and the chamber graph of every residue, is connected: by induction on rank, chambers containing a given element are connected; a path in the incidence graph connects these families, using chambers containing each consecutive incident pair. The rank-one and rank-two cases start the induction. Thus a chamber stabilizer is trivial. The same argument inside the residue of the base flag of types $J$ proves (\ref{eq:parabolic}). It also proves that the full type-preserving automorphism group is the displayed group whenever a faithful group already acts transitively on chambers.

For a second regular hypertope $(A,(A_x)_{x\in S})$, put $A(T)=\langle a_x:x\in T\rangle$ for $T\subseteq S$. Formula (\ref{eq:parabolic}) gives
\begin{equation}\label{eq:ip}
A(T)\cap A(T')=A(T\cap T').
\end{equation}
We will use both this intersection property and flag-transitivity. They are different requirements: the intersection property alone is not a general criterion for a nonstring coset system to be a geometry.

\begin{lemma}[Coset intersection]\label{lem:helly}
If $B$ is flag-transitive on its coset geometry, every finite pairwise-intersecting family of cosets of standard parabolic subgroups has nonempty total intersection. In particular, for $H,K,J\subseteq I$,
\[
B_HB_J\cap B_KB_J=B_{H\cup K}B_J.
\]
\end{lemma}
\begin{proof}
Expand a coset $gB_H$ as $\bigcap_{i\in H}gB_i$. Pairwise intersection of the original cosets implies pairwise intersection of all the resulting maximal-parabolic cosets. Intersecting cosets of the same type are equal. They therefore form a flag, which flag-transitivity takes to a standard flag; its cosets have a common group element. For the last assertion, if $g\in B_HB_J\cap B_KB_J$, the three cosets $B_H,B_K,gB_J$ intersect pairwise. A point of their common intersection belongs to $B_{H\cup K}$ and to $gB_J$. The reverse inclusion follows from subgroup containment.
\end{proof}

\begin{theorem}[Element orbits]\label{thm:elements}
Fix $j\in I$ and $x_0=B_j\in X=B/B_j$. If $B$ is flag-transitive, then
\begin{equation}\label{eq:orbit}
(B_Hx_0)\cap(B_Kx_0)=B_{H\cup K}x_0\qquad(H,K\subseteq I).
\end{equation}
Thus the complementary notation $O(M)=B_{I\setminus M}x_0$ satisfies $O(M)\cap O(N)=O(M\cap N)$.
\end{theorem}
\begin{proof}
The inverse image of $B_Hx_0$ under $B\to B/B_j$ is $B_HB_j$. Apply Lemma~\ref{lem:helly} and take images. Finally $(I\setminus M)\cup(I\setminus N)=I\setminus(M\cap N)$.
\end{proof}

\input{general_twisting.tex}

\section{Invariant quotients and orbital presentations}
\begin{lemma}[Quotients]\label{lem:quotient}
If $M\leq A$ is invariant under $B$, its normal closure $N$ in $A$ is invariant. The action descends to $A/N$ and permutes the images of the standard maximal parabolics with the same indexed type action. If their coset geometry is a regular hypertope, the corresponding twisting is again a regular hypertope.
\end{lemma}
\begin{proof}
The normal closure is generated by $a m a^{-1}$ with $a\in A,m\in M$. An automorphism from $B$ carries each such generator to another of the same form, and its inverse gives equality. The quotient action and parabolic permutation follow. Orbit identity (\ref{eq:general-orbits}) concerns only the indexed types and is unchanged.

For completeness, the regular-geometry assumption can be interpreted without silently assuming a faithful quotient action. Put $\bar A=A/N$, $C_x=A_xN/N$, and $C_E=\bigcap_{x\in E}C_x$, with $C_\varnothing=\bar A$. The source's quotient twisting uses intersections of these maximal parabolics:
\begin{equation}\label{eq:quotientparabolics}
\bar G_i=C_{S\setminus\bigcup_L B_ix_L}\rtimes B_i\quad(i\in I),
\qquad \bar G_L=C_L\rtimes B\quad(L\in\mathcal K).
\end{equation}
These are not defined by substituting the images of generated-support subgroups into (\ref{eq:regular-support}).

Each image of $a_x$ fixes the base flag with type $x$ deleted. It cannot belong to $C_x$: otherwise $C_x$, already containing the other generators, would be all of $\bar A$, giving only one element of type $x$, contrary to thinness. It therefore performs the unique adjacency of that type. Chamber-graph connectivity makes $\bar A$ flag-transitive. Its chamber stabilizer $H=C_S$ fixes every chamber by thinness, so is its normal action kernel. It is $B$-invariant because it is the intersection of the permuted $C_x$.

Every $C_E$ contains $H$, and intersections commute with passage to $\bar A/H$. In that faithful regular model formula (\ref{eq:parabolic}) therefore gives
\[
C_E/H=\bigcap_{x\in E}(C_x/H)
=\langle a_xH:x\in S\setminus E\rangle.
\]
Consequently every subgroup in (\ref{eq:quotientparabolics}) is the full inverse image of the corresponding subgroup from Theorem~\ref{thm:twist} for $\bar A/H$. Apply that theorem. Passage through the normal subgroup $H\rtimes1$ preserves all cosets and their intersections, so the original quotient twisting has the same regular geometry. For a transitive type orbit $L=S$, its new parabolic is $H\rtimes B$, not necessarily $B$. The quotient group $\bar A$ itself is not asserted to be the faithful automorphism group.
\end{proof}

\begin{proposition}[Orbital presentation]\label{prop:orbital}
Let $B$ act transitively on a finite set $X$ with base point $x_0$ and stabilizer $P$. Give each orbit of unordered distinct pairs a Coxeter label. Let $Y$ be a finite set fixed by $B$, with fixed labels within $Y$ and one label $m_y$ between $y$ and every element of $X$. Put $A=W(Y\sqcup X)$. Then $A\rtimes B$ has the presentation consisting of a presentation for $B$, one involution $a=a_{x_0}$, the involutions $c_y$ for $y\in Y$, and:
\begin{enumerate}
\item $[a,p]=1$ for generators $p$ of $P$, and $[c_y,b]=1$ for generators $b$ of $B$;
\item the prescribed relations among the $c_y$, and $(c_ya)^{m_y}=1$ when $m_y<\infty$;
\item $(a\tau a\tau^{-1})^{m}=1$ for one representative $\{x_0,\tau x_0\}$ of each unordered pair orbit of finite label $m$.
\end{enumerate}
The unordered orbits are indexed by the inverse-paired double cosets $P\tau P\cup P\tau^{-1}P$ with $\tau\notin P$.
\end{proposition}
\begin{proof}
In the proposed group define $a_{bx_0}=bab^{-1}$. The stabilizer relations make this independent of $b$. Double-coset multiplication changes the representative without changing its pair orbit, and inversion interchanges its endpoints. Conversely, a transporter of two unordered pairs either fixes the selected first endpoint or interchanges the endpoints, giving exactly the two double cosets displayed. Hence conjugates of the chosen relations give every Coxeter relation on $X$. Conjugating the relations with $c_y$, which is fixed by $B$, gives all the mixed Coxeter relations. These elements therefore define a homomorphism from $W(Y\sqcup X)\rtimes B$ onto the proposed group. The reverse map takes the proposed generators to their named elements in the semidirect product; all the relators hold there. Both composites fix the generating set, proving the isomorphism. This is a presentation of the full group, not merely a finite image.
\end{proof}

\begin{lemma}[Exact doubled label]\label{lem:double}
In $A\rtimes B$ let $a=a_x$ and let $b\in B$ be an involution with $bx\ne x$. If $a_xa_{bx}$ has order $m$ (possibly infinite), then $ab$ has order $2m$, with $2\infty=\infty$.
\end{lemma}
\begin{proof}
We have $(ab)^2=a_xa_{bx}$. Projection to $B$ sends $ab$ to the nontrivial involution $b$, so any power of $ab$ equal to $1$ has even exponent. For an even exponent $2k$ it is $1$ precisely when $(a_xa_{bx})^k=1$. This proves exactness in both the finite and infinite cases.
\end{proof}

\input{module_shears.tex}

\section{Attribution and retained facts}
The constructions and target results are due to Piedade and Tranchida [1,2]. The orbital presentation is the same argument used for the matching extensions in our earlier common paper, with the fixed vertices included explicitly. The contribution of this proof edition is a common general semidirect theorem, exact panel formulas and explicit adapters, not a priority claim for generalized twisting. The marked-module and cocycle arguments expose the common proof of the signed families; their application to the odd-rank case is stated separately from known recoveries. The classical Coxeter facts used by the application are that standard parabolic subgroups embed with their indicated Coxeter presentations and that the standard coset geometry is a regular hypertope. The usual string C-group--regular-polytope correspondence and the named classical twisting constructions are retained as stated in the application. The companion Lean file verifies only the declared coset, support, and order identities; it is not an end-to-end geometry formalization.
\begin{thebibliography}{9}
\bibitem{PTj} C. A. Piedade and P. Tranchida, \emph{Coset geometries acting on their elements: The $j$-diagonals twisting}, arXiv:2608.23223v1, 2026.
\bibitem{PTw} C. A. Piedade and P. Tranchida, \emph{The geometry of wreath and semi-direct products}, arXiv:2605.04646v2, 2026, Definition 3.10 and Corollary 3.17.
\bibitem{FP} M. E. Fernandes and C. A. Piedade, \emph{Regular polytopes of rank $n/2$ for transitive groups of degree $n$}, arXiv:2407.16003v2, 2025; J. Algebra 689 (2026), 256--296. The diagram correspondences use Tables 6 and 9 of the bound arXiv edition.
\end{thebibliography}
\end{document}
